% SOURCE_AUTHORITY: sga5_fr_workpass.tex, lines 8631--8647 (LF-normalized unit).
% SOURCE_SHA256: 791F4EFFC5E02832D5D77ED03518C8156D6F07E4C8238B03545DB93D883FBB28
\subsubsection*{1.4.4}
Se \(A\) é uma \(\mathbb Z_\ell\)-álgebra comutativa finita e \(F\) e \(G\) são dois \(A\)-feixes construtíveis, o sistema projetivo
\[
F\otimes_A G=(F_n\otimes_{A_X} G_n)_{n\in\mathbb N},
\]
com os morfismos de transição evidentes, é um \(\mathbb Z_\ell\)-feixe. Além disso, \(A\) atua de maneira evidente sobre cada uma de suas componentes, o que permite dotá-lo naturalmente de uma estrutura de \(A\)-feixe. Define-se assim um bifuntor exato à direita, chamado produto tensorial,
\[
A\text{-}\fc(X)\times A\text{-}\fc(X)\longrightarrow A\text{-}\fc(X).
\]

Sejam agora \(F\) e \(G\) dois \(A\)-feixes e suponhamos que \(F\) é localmente livre, isto é, que, para todo inteiro \(n\), \(F_n\) é um \((A/\ell^{\,n+1}A)_X\)-módulo localmente livre. Então, como em (1.3.3), vê-se que o sistema projetivo
\[
\cHom_A(F,G)=\bigl(\cHom_{A_n}(F_n,G_n)\bigr)_{n\in\mathbb N},
\]
com os morfismos de transição evidentes, é um \(\mathbb Z_\ell\)-feixe. Além disso, como \(A\) age naturalmente sobre suas componentes, ele está canonicamente dotado de uma estrutura de \(A\)-feixe.

Tal como em (1.3.4), define-se a noção de \(A\)-feixe invertível e, se \(L\) é tal \(A\)-feixe, dá sentido à expressão \(L^{\otimes r}\) para todo inteiro natural \(r\). O n.o 1.3.4 adapta-se sem dificuldade a esta situação.
